% Source: https://artofproblemsolving.com/community/c3h2784571_2021_abmc_team_rounds_58__actonboxborough_math_competition
% Stable identity: aops-topic-2784571-q19
Write $x=gu$, $y=gv$ with $\gcd(u,v)=1$. Then $\operatorname{lcm}(2u,v)=2uv$ if $v$ is odd and $uv$ if $v$ is even, so $\ell=e\,guv$ with $e=2$ or $1$ accordingly. Putting $K=uv$, the equation becomes
\[g(gK+3+7eK)=168,\qquad K(g+7e)=\dfrac{168}g-3.\]
Checking the divisors $g$ of $168$:
\begin{itemize}
\item $g=1$: $K(1+7e)=165$. With $e=2$, $K=11$; with $e=1$, $8K=165$ fails. So $uv=11$ with $v$ odd: $(x,y)=(1,11)$ or $(11,1)$.
\item $g=7$: $K(7+7e)=21$. With $e=2$, $K=1$, so $(x,y)=(7,7)$.
\item $g=2,3,4,6,8,12,14,21,24,28,42,56,84,168$: the right side is $81,53,39,25,18,11,9,5,4,3,1,0$ or negative, and neither $K(g+14)$ nor $K(g+7)$ matches with $K\ge1$ (for $g=2$, $9K=81$ gives $uv=9$ with $v$ even, which is impossible).
\end{itemize}
The candidates give $2x+y=13$, $23$ and $21$, so the maximum is $\boxed{23}$, from $(x,y)=(11,1)$: $11+3+7\cdot22=168$.
